\topictitle{\tr{Funcția de gradul doi}{The quadratic function}{Die quadratische Funktion}}

\[
f(x)=ax^2+bx+c,\qquad \Delta=b^2-4ac
\]

\subsection{\tr{Tabel de semn}{Sign table}{Vorzeichentabelle}}
\tr{Pentru}{For}{Für}\ $\Delta>0$, $x_1<x_2$:
\[
\begin{array}{c|ccccccc}
x & -\infty & & x_1 & & x_2 & & +\infty\\ \hline
f(x) & & \mt{semnul lui $a$}{sign of $a$}{Vorzeichen von $a$} & 0 &
 \mt{semn contrar lui $a$}{opposite sign of $a$}{Vorzeichen von $a$ umgekehrt} & 0 &
 \mt{semnul lui $a$}{sign of $a$}{Vorzeichen von $a$} &
\end{array}
\]

\subsection{\tr{Parabola (vârful)}{The parabola (vertex)}{Die Parabel (Scheitelpunkt)}}
\[
V\Bigl(-\frac{b}{2a},\,-\frac{\Delta}{4a}\Bigr)=V(x_v,y_v)
\]

\subsection{\tr{Valoarea minimă / maximă}{Minimum / maximum value}{Minimum / Maximum}}
\[
\begin{aligned}
a>0&:\ \mt{valoarea minimă}{minimum value}{Minimum}\ -\frac{\Delta}{4a}\quad\mt{pentru}{at}{bei}\ x=-\frac{b}{2a}\\[2pt]
a<0&:\ \mt{valoarea maximă}{maximum value}{Maximum}\ -\frac{\Delta}{4a}\quad\mt{pentru}{at}{bei}\ x=-\frac{b}{2a}
\end{aligned}
\]

\subsection{\tr{Monotonia}{Monotonicity}{Monotonie}}
\[
a>0:\ \begin{cases}\text{\tr{strict descrescătoare pe}{strictly decreasing on}{streng monoton fallend auf}}\ \bigl(-\infty,-\tfrac{b}{2a}\bigr]\\[2pt]
\text{\tr{strict crescătoare pe}{strictly increasing on}{streng monoton steigend auf}}\ \bigl[-\tfrac{b}{2a},\infty\bigr)\end{cases}
\qquad a<0:\ \mt{invers}{reversed}{umgekehrt}
\]

\begin{important}
\tr{Ca funcția să se anuleze într-un singur punct: $\Delta=0$.}{For the function to vanish at exactly one point: $\Delta=0$.}{Damit die Funktion genau eine Nullstelle hat: $\Delta=0$.}
\end{important}

\subsection{\tr{Familii de parabole (cu parametru)}{Families of parabolas (with a parameter)}{Parabelscharen (mit Parameter)}}
\begin{itemize}
\item \tr{Vârfurile parabolelor se află pe:}{The vertices lie on:}{Die Scheitelpunkte liegen auf:}
  \begin{itemize}[label=$\blacktriangleright$]
  \item \tr{o dreaptă ($y_v$ este funcție de gradul 1 de $x_v$)}{a line ($y_v$ is a degree-1 function of $x_v$)}{einer Geraden ($y_v$ ist eine Funktion 1.\ Grades von $x_v$)}
  \item \tr{o parabolă ($y_v$ este funcție de gradul 2 de $x_v$)}{a parabola ($y_v$ is a degree-2 function of $x_v$)}{einer Parabel ($y_v$ ist eine Funktion 2.\ Grades von $x_v$)}
  \item \tr{prima bisectoare:}{the first bisector:}{der ersten Winkelhalbierenden:}\ $x_v=y_v$
  \item \tr{a doua bisectoare:}{the second bisector:}{der zweiten Winkelhalbierenden:}\ $x_v=-y_v$
  \item \tr{dreapta $\dots$ (înlocuim $x$ și $y$ cu $x_v$ și $y_v$)}{a line $\dots$ (replace $x$ and $y$ by $x_v$ and $y_v$)}{einer Geraden $\dots$ ($x$ und $y$ durch $x_v$ und $y_v$ ersetzen)}
  \end{itemize}
\item \tr{Curba pe care se găsesc vârfurile: exprimăm parametrul în funcție de $x_v$ și îl înlocuim în ecuația lui $y_v$.}
        {The curve containing the vertices: express the parameter through $x_v$ and substitute it into the equation of $y_v$.}
        {Ortskurve der Scheitelpunkte: den Parameter durch $x_v$ ausdrücken und in die Gleichung für $y_v$ einsetzen.}
\item \tr{Parabolele trec printr-un punct fix (cel puțin unul): scoatem parametrul factor comun și egalăm cu $0$ coeficienții acestuia.}
        {The parabolas pass through a fixed point (at least one): factor out the parameter and set its coefficients equal to $0$.}
        {Die Parabeln gehen durch einen Fixpunkt (mindestens einen): den Parameter ausklammern und seine Koeffizienten gleich $0$ setzen.}
\end{itemize}

\subsection{\tr{Parabolele a două funcții}{Parabolas of two functions}{Parabeln zweier Funktionen}}
\begin{itemize}[label=$\blacktriangleright$]
\item \tr{se intersectează în două puncte distincte:}{intersect in two distinct points:}{schneiden sich in zwei verschiedenen Punkten:}
  $\begin{cases}f(x)=g(x)\\ \Delta_{f-g}>0\end{cases}$
\item \tr{au un singur punct comun:}{have exactly one common point:}{haben genau einen gemeinsamen Punkt:}
  $\begin{cases}f(x)=g(x)\\ \Delta_{f-g}=0\end{cases}$
\end{itemize}

\subsection{\tr{Grafic}{Graph}{Graph}}
\begin{center}
\begin{tikzpicture}[scale=1]
  \begin{scope}
    \draw[->,gray] (-1.3,0)--(1.4,0); \draw[->,gray] (0,-0.2)--(0,1.5);
    \draw[thick,accent,domain=-1.1:1.1,smooth,samples=40] plot (\x,{1.2*\x*\x});
    \node at (0,-0.6) {$a>0$};
  \end{scope}
  \begin{scope}[xshift=5cm]
    \draw[->,gray] (-1.3,0)--(1.4,0); \draw[->,gray] (0,-1.5)--(0,0.2);
    \draw[thick,accent,domain=-1.1:1.1,smooth,samples=40] plot (\x,{-1.2*\x*\x});
    \node at (0,0.6) {$a<0$};
  \end{scope}
\end{tikzpicture}
\end{center}

\subsection{\tr{Mărginirea}{Boundedness}{Beschränktheit}}
\begin{itemize}[label=$\blacktriangleright$]
\item \tr{Inferior}{Bounded below}{Nach unten beschränkt}\ $\imp\exists\,m\in\R:\ m\le f(x)\ \ \forall x$
\item \tr{Superior}{Bounded above}{Nach oben beschränkt}\ $\imp\exists\,M\in\R:\ M\ge f(x)\ \ \forall x$
\end{itemize}
